% Local charged-configuration estimate used in the trace response bound.
\section{The local charged-configuration estimate}
\label{sec:ex-lattice-trace-local-charge}

We isolate the finite-volume algebra that turns a charged configuration
into a bounded field and a priced response. This is the local step used
before the separation-limit passage. Its quantitative premises are stated
as component bounds, so the proof does not hide the estimate inside a
single named receipt.

For $t>e$, put
\[
 W_a(t)=\exp\bigl(a(\log\log t)^2\bigr),
 \qquad c(t)=t^{-30}.
\]
Then $W_a(t)W_b(t)=W_{a+b}(t)$, $W_{-1}(t)\to0$, and $c(t)\to0$ as
$t\to\infty$.

Let $I$ be a finite set of lattice directions, $Y$ a finite-dimensional
real normed coefficient space, and let $\ell:Y^I\to\mathbb R$ be a
continuous linear flux functional. Let $s:Y^I\to[0,\infty)$ be the source coefficient size, satisfying
$\|A\|_\infty\le s(A)$ for every coefficient array $A$. Write
$\sigma=\theta^{-2}$ with $\theta>0$ for the local normalization. The
nonlinear local coefficient map has already been evaluated in the arrays
$B',K',D'\in Y^I$; these are outputs of that local construction, not images
under the core linear projection. Let $X,W,E,L\in Y^I$. Assume the exact
charged-configuration decomposition
\begin{equation}
 X=W+B'-K'+D'+E+L,
 \label{eq:ex-lattice-charge-decomposition}
\end{equation}
and the following local coefficient and flux receipts at scale $t$:
\begin{align}
 s(X)&\le\theta^2 W_{1/6}(t),
 \label{eq:ex-lattice-charge-size-input}\\
 \ell(\sigma K')&=0,
 &\ell(\sigma B')&=h/12,
 &|h-m|&\le c(t),
 \label{eq:ex-lattice-charge-price-identities}\\
 s(W)&\le\theta^2 W_{-1}(t),
 &s(D')&\le\theta^2 W_{-1}(t),
 \label{eq:ex-lattice-charge-error-one}\\
 s(E)&\le\theta^2 W_{-1}(t),
 &s(L)&\le\theta^2 W_{-1}(t).
 \label{eq:ex-lattice-charge-error-two}
\end{align}
Here $h$ is the finite-volume charge extracted from the distinguished
component, and $m$ is its target price. These are coefficient-size receipts
after the local map has been evaluated. No norm contraction of the nonlinear
map is assumed.

\begin{proposition}[Local field and flux bounds]
\label{prop:ex-lattice-charge-bounds}
Under the displayed assumptions,
\[
 \|\sigma X\|_\infty\le W_{1/6}(t),
 \qquad
 \left|\ell(\sigma X)-\frac m{12}\right|
 \le4\|\ell\|W_{-1}(t)+\frac{c(t)}{12}.
\]
\end{proposition}
\begin{proof}
The field bound follows from the size domination and
\eqref{eq:ex-lattice-charge-size-input}:
\[
 \|\sigma X\|_\infty\le\sigma s(X)=\theta^{-2}s(X)\le W_{1/6}(t).
\]
For the price, apply $\ell$ to the decomposition in
\eqref{eq:ex-lattice-charge-decomposition}. The $K'$ term vanishes and the
$B'$ term is $h/12$. Therefore
\[
 \ell(\sigma X)-\frac m{12}
 =\ell(\sigma W)+\ell(\sigma D')
   +\ell(\sigma E)+\ell(\sigma L)+\frac{h-m}{12}.
\]
For any array $A$, continuity gives
$|\ell(\sigma A)|\le\|\ell\|\sigma\|A\|_\infty\le\|\ell\|\theta^{-2}s(A)$. The triangle
inequality and
\eqref{eq:ex-lattice-charge-price-identities}--
\eqref{eq:ex-lattice-charge-error-two} give the claimed bound. There are
four error terms, so the coefficient is exactly four; no direction-count
factor is introduced because $s$ bounds the full array norm.
\end{proof}

In the selected trace response, $m/12$ is the target flux $-\kappa$.
Consequently the local bound proves a response-size estimate tending to
zero after division by the arch scale and an unweighted flux limit, provided
$t\to\infty$ along the selected configurations and the displayed component
receipts hold. The remaining passage from typical finite-volume
configurations to their separated limit is an ordinary continuity
argument, as recorded in the weighted-estimate section. The component
receipts themselves are separate finite-volume producer statements; this
proposition does not prove them from the charged-configuration hypotheses.
